%# -*- coding:utf-8 -*-
%% start of file `template_en.tex'.
%% Copyright 2006-1008 Xavier Danaux (xdanaux@gmail.com).
%
% This work may be distributed and/or modified under the
% conditions of the LaTeX Project Public License version 1.3c,
% available at http://www.latex-project.org/lppl/.


\documentclass[11pt,a4paper]{moderncv}

\usepackage{fontspec,xunicode}
\usepackage[slantfont,boldfont]{xeCJK}
\usepackage{xcolor}  % 颜色支持


% \setmainfont{Tahoma}
\setmainfont{Times New Roman}  % 缺省英文字体；serif 为衬线字体，sans serif 为无衬线字体

% macOS 字体设置
\setCJKmainfont[ItalicFont={Kai}, BoldFont={Hei}]{STSong}  % 衬线字体 缺省中文字体
\setCJKsansfont{STSong}
\setCJKmonofont{STFangsong}  % 中文等宽字体

%-----------------------xeCJK 下设置中文字体（macOS）------------------------------%
\setCJKfamilyfont{song}{Songti SC}  % 宋体 song
\newcommand{\song}{\CJKfamily{song}}
\setCJKfamilyfont{fs}{STFangsong}  % 仿宋 fs
\newcommand{\fs}{\CJKfamily{fs}}
\setCJKfamilyfont{yh}{PingFang SC}  % 苹方 yh
\newcommand{\yh}{\CJKfamily{yh}}
\setCJKfamilyfont{hei}{Heiti SC}  % 黑体 hei
\newcommand{\hei}{\CJKfamily{hei}}
\setCJKfamilyfont{hwxh}{STHeiti}  % 华文黑体 hwxh
\newcommand{\hwxh}{\CJKfamily{hwxh}}
\setCJKfamilyfont{kai}{Kaiti SC}  % 楷体 kai
\newcommand{\kai}{\CJKfamily{kai}}


%------------------------------设置字体大小------------------------%
\newcommand{\chuhao}{\fontsize{42pt}{\baselineskip}\selectfont}  % 初号
\newcommand{\xiaochuhao}{\fontsize{36pt}{\baselineskip}\selectfont}  % 小初号
\newcommand{\yihao}{\fontsize{28pt}{\baselineskip}\selectfont}  % 一号
\newcommand{\erhao}{\fontsize{21pt}{\baselineskip}\selectfont}  % 二号
\newcommand{\xiaoerhao}{\fontsize{18pt}{\baselineskip}\selectfont}  % 小二号
\newcommand{\sanhao}{\fontsize{15.75pt}{\baselineskip}\selectfont}  % 三号
\newcommand{\sihao}{\fontsize{14pt}{\baselineskip}\selectfont}  % 四号
\newcommand{\xiaosihao}{\fontsize{12pt}{\baselineskip}\selectfont}  % 小四号
\newcommand{\wuhao}{\fontsize{10.5pt}{\baselineskip}\selectfont}  % 五号
\newcommand{\subwuhao}{\fontsize{10pt}{\baselineskip}\selectfont}  % 次五号
\newcommand{\xiaowuhao}{\fontsize{9pt}{\baselineskip}\selectfont}  % 小五号
\newcommand{\liuhao}{\fontsize{7.875pt}{\baselineskip}\selectfont}  % 六号
\newcommand{\qihao}{\fontsize{5.25pt}{\baselineskip}\selectfont}  % 七号


% \usepackage{fontawesome}
% \setCJKmainfont[BoldFont={WenQuanYi Micro Hei/Bold}]{WenQuanYi Micro Hei}
% \defaultfontfeatures{Mapping=tex-text}
% \XeTeXlinebreaklocale "zh"
% \XeTeXlinebreakskip = 0pt plus 1pt minus 0.1pt
% moderncv 主题
\moderncvtheme[blue]{classic}  % 颜色可选 'blue'（默认）、'orange'、'red'、'green'、'grey'，或 'roman'（使用罗马字体）
% \moderncvtheme[green]{classic}  % 同上
% \moderncvtheme[blue,roman]{hht}
% 字符编码



% 调整页面边距
\usepackage[scale=0.9]{geometry}
% \setlength{\hintscolumnwidth}{3cm}  % 如需调整左侧日期栏宽度可取消注释
% \AtBeginDocument{\setlength{\maketitlenamewidth}{6cm}}  % 仅 classic 主题有效，用于调整姓名占位宽度（为右侧联系信息留出更多空间）
\AtBeginDocument{\recomputelengths}  % 修改页面尺寸相关参数后必须调用

% 个人信息
% 中文简历不区分姓与名，直接把完整姓名填入 \firstname 即可，
% \familyname 留空。moderncv 内部需要这两个字段同时存在。
\firstname{张三}                 % 完整姓名
\familyname{}                    % 留空（moderncv 内部需要此字段）
\title{SanZhang}                 % 头衔 / 拼音名，不需要可注释本行
% \address{杭州}{}                % 地址，不需要可注释本行
\address{1990/11/11}{}           % 此处用于填出生日期，不需要可注释本行
\mobile{18888888888}             % 手机号，不需要可注释本行
% \fax{fax (optional)}           % 传真，不需要可注释本行
\email{me@resume.com}            % 邮箱，不需要可注释本行
\homepage{Blog: http://geekplux.com}  % 主页 / 博客，不需要可注释本行
\extrainfo{%
  GitHub: \url{https://github.com/geekplux} \\
  LinkedIn: \url{https://cn.linkedin.com/in/xxx} \\
  WeChat: xxxx \\
  QQ: 123456
}

\photo[64pt]{avatar.png}  % 头像；'64pt' 为缩放后的高度，第二个参数是图片文件名。不需要头像可注释本行
% \quote{China\TeX 您的LaTeX乐园，TeX\&\LaTeX 王国}  % 个性签名 / 引言，不需要可注释本行

% \nopagenumbers{}  % 当简历超过一页时取消注释可关闭自动页码


%----------------------------------------------------------------------------------
%            正文内容
%----------------------------------------------------------------------------------
\begin{document}
\maketitle
\vspace*{-14mm}

\section{工作经历}
\cventry{17.01-17.12}{不长不短的公司名称}{可爱的项目}{https://keaidexiangmu.com}{}{工作中，
  我负责了 xxx 的开发和维护，运用 yyy 技术解决了 zzz 的重大问题。积极参与开源社
  区贡献}
\cventry{16.01-16.12}{好长的公司名称}{有趣的项
  目}{http://www.youqudexiangmu.com}{}{工作中，我按照领导的要求编程，做出了让领
  导满意的作品，为公司做出了贡献。我觉得我还是凑字数吧，编不下去了}
\cventry{15.01-15.12}{不知道叫什么的公司名称}{不可告人的项
  目}{http://bukegaoren.com}{}{独立编写了根本编不下去的项目简介，可能还是凑字数
  比较好，来凑字数吧来凑字数吧来凑字数吧。出色的完成了凑字数的工作，并获得了最佳
  凑字数员工奖}
\section{教育经历}
\cventry{11.09-15.06}{本科}{尼姑庵大学}{计算机科学与技术}{}{}  % 第 3-6 个参数可选
\cvlistitem{最快编程大师一等奖}
\cvlistitem{最强编程大师金奖}
\cvlistitem{第 x 届「编程杯」gayhub 赛区一等奖}
\cvlistitem{国家奖学金/三好学生/学生会主席/\emph{获得女朋友一个}}
\cventry{15.09-18.06}{硕士}{和尚庙大学}{软件工程}{实验室 XXX 导师 XXX}{主要研究了
  人工智能，图形学，编译原理，机械键盘的拆装，快递包装的暴力拆解，颈椎与视觉保
  养，抹平小腹，治疗腰椎间盘突出}  % 第 3-6 个参数可选
\section{项目}
\subsection{科研项目}
\cvline{FFF-TNT}{建立深度学习模型，通过训练单身狗愤怒指数，预测 TNT 威力值，达到
  最后 FFF 全部的效果}
\cvline{NASA}{我又编不下去了，还是凑个字数，要不然你们将来替换也麻烦}
\subsection{开源项目}
\cvline{markvis}{在 markdown 中直接生成可视化图表的插件
  \emph{https://markvis.js.org} \textbf{GitHub 1000 stars}}
\cvline{netjsongraph\-.js}{用力导向图可视化出无线路由图谱数据 \emph{https://github.com/netjson/netjsongraph.js}}
\cvline{typing}{Hexo 静态博客主题 \emph{https://github.com/geekplux/hexo-theme-typing}}
\cvline{UnityVis}{Unity 中的基本可视化图表 \emph{https://github.com/geekplux/Basic-Visualization-in-Unity}}
\section{技能}
\cvline{\textbf{前端}}{熟练掌握前端该会的东西，熟悉前端该熟悉的东西，了解前端该
  了解的原理，熟练使用前端该熟练使用的各类工具}
\cvline{\textbf{后端}}{熟练掌握后端该会的东西，熟悉后端该熟悉的东西，了解后端该了解的原理}
\cvline{\textbf{数据}}{掌握处理数据，掌握数据库}
\cvline{\textbf{其他}}{熟练使用 Git/Vim/Emacs/Makefile}

\section{Publications}
\cvline{已录用}{张三，李四，王麻子 基于 latex 的简历凑字数研究[C]// CVChina. 2017.}

% 以下为可选章节示例，需要时取消注释即可启用
% \subsection{职业经历}
% \cventry{年份--年份}{职位}{雇主}{城市}{}{描述}  % 第 3-6 个参数可选
% \cventry{年份--年份}{职位}{雇主}{城市}{}{描述}  % 第 3-6 个参数可选
% \subsection{其他经历}
% \cventry{年份--年份}{职位}{雇主}{城市}{}{描述第一行\newline{}描述第二行}  % 第 3-6 个参数可选

% \section{语言}
% \cvlanguage{语言 1}{熟练程度}{备注}
% \cvlanguage{语言 2}{熟练程度}{备注}
% \cvlanguage{语言 3}{熟练程度}{备注}

% \section{计算机技能}
% \cvcomputer{类别 1}{XXX, YYY, ZZZ}{类别 4}{XXX, YYY, ZZZ}
% \cvcomputer{类别 2}{XXX, YYY, ZZZ}{类别 5}{XXX, YYY, ZZZ}
% \cvcomputer{类别 3}{XXX, YYY, ZZZ}{类别 6}{XXX, YYY, ZZZ}

% \section{兴趣爱好}
% \cvline{篮球}{\small 体力与技巧}
% \cvline{爱好 2}{\small 描述}
% \cvline{爱好 3}{\small 描述}

% \renewcommand{\listitemsymbol}{-}  % 修改列表项符号

% \section{额外章节 1}
% \cvlistitem{项目 1}
% \cvlistitem{项目 2}
% \cvlistitem[+]{项目 3}  % 自定义列表符号
% XeLaTeX 支持任意 macOS 字体，详见上文 \setromanfont 命令
% XeLaTeX 输入支持完整 Unicode，可直接输入 Unicode 字符

% 以下两行用于告诉 TeXShop 使用 xelatex 编译，并使用 UTF-8 编码读写源文件
% !TEX TS-program = xelatex
% !TEX encoding = UTF-8 Unicode

% \section{额外章节 2}
% \cvlistdoubleitem[\Neutral]{项目 1}{项目 4}
% \cvlistdoubleitem[\Neutral]{项目 2}{项目 5}
% \cvlistdoubleitem[\Neutral]{项目 3}{}

%% 从 BibTeX 文件导入论文列表
% \nocite{*}
% \bibliographystyle{plain}
% \bibliography{publications}  % 'publications' 是 BibTeX 文件名

% \begin{thebibliography}{}
% \bibitem[]{} 移动增强现实可视化综述[C]. ChinaVis 2017.
% \end{thebibliography}


\end{document}


%% end of file `template_en.tex'.

%%% Local Variables:
%%% mode: latex
%%% TeX-command-extra-options: "-shell-escape"
%%% TeX-master: t
%%% TeX-engine: xetex
%%% End:
